\begin{abstract}
A benchmark score tells us how well a driving policy performed on logged scenes, much as a return
figure tells us how a fund performed. Before deploying either, one wants a risk profile. How much does
it put at stake? Does it manage risk when risk appears? Does it grow more careful or more reckless as
conditions heat up? How much does noise move it? We propose a check-up that reads such a profile from
a few hundred logged frames, without a simulator and without retraining. Each exam edits the frames
counterfactually, removing a pedestrian from the ego corridor or re-rendering the scene at night, and
pairs the edit with a control. An independent detector verifies every edit. The central exam is a
hazard-sensitivity curve: how much farther the whole planned trajectory stays from the pedestrian's
logged future because the policy can see the pedestrian, as a function of hazard. We examine six
end-to-end policies. Their exposure varies fourfold (\NumExpLo--\NumExpHi{} times the human's
distance), yet their hazard sensitivity is \NumHSlo--\NumHShi{} on a $[-1,1]$ scale and does not rise
with hazard. Collisions follow exposure, not perception. A night rendering moves every policy more
than the pedestrian does ($\CFR=\NumCFRlo$--$\NumCFRhi$), and three of them plan up to
\NumNightDD\% faster in the dark. Fitting estimator spread against sample
size shows that exposure, over-reaction and the lighting verdict are settled within about 50 frames,
while hazard sensitivity needs hundreds. In a pre-registered test, the orderings and verdicts
diagnosed on Boston frames carry over to Singapore, but the point values do not. Against the standard
evaluations, the check-up agrees where they measure the same thing and shows what they miss: nuScenes
open-loop metrics do not change when the pedestrian is removed, and NAVSIM's ranking reverses on
scenes with a pedestrian nearby.
\end{abstract}
